\documentclass[10pt,a4paper]{amsart}
\usepackage[margin=19mm]{geometry}
\usepackage{amsmath,amssymb,mathtools}
\usepackage[scr=rsfs,cal=euler]{mathalfa}
\usepackage{xfrac}
\usepackage[unicode,hidelinks,bookmarks=false]{hyperref}
\newcommand{\ie}{i.e.\ }
\newcommand{\cf}{cf.\ }
\newcommand{\resp}{respectively\ }
\newcommand{\Subsection}{\subsection}
\providecommand{\nobreakdash}{\nobreak\hbox{-}}
\setlength{\emergencystretch}{2em}
\multlinegap0pt
\usepackage{mathtools}
\usepackage{xcolor}
\usepackage{amssymb}
\usepackage{bm}
\usepackage[shortlabels]{enumitem}
\usepackage{algorithm}\usepackage{algpseudocode}
\newtheorem{proposition}{Proposition}
\newcommand{\RR}{\mathbb{R}}
\DeclareMathOperator{\dist}{dist}
\title{Retraction based regression methods on Riemannian manifolds}
\author{Estefan\'ia Loayza-Romero, Benedikt Sibum, Kathrin Welker}
\date{Source: arXiv:2606.02116v1 (CC BY 4.0)}
\begin{document}
\maketitle

\noindent\emph{Source and licence.} Retraction based regression methods on Riemannian manifolds. Version arXiv:2606.02116v1 (CC BY 4.0), distributed by arXiv under the Creative Commons Attribution 4.0 International licence, http://creativecommons.org/licenses/by/4.0/, as recorded in the arXiv licence metadata for this exact version and shown on the abstract page https://arxiv.org/abs/2606.02116v1. Full licence notice: \texttt{licenses/arxiv-2606.02116-CC-BY-4.0.txt}.\par\smallskip

\noindent\textit{Editorial scope.} Counted unit: the article's proposition prop:well-posed (source0.tex lines 219--229). The article's complete original proof of it is delivered with it (source0.tex lines 230--232, one \texttt{proof} environment of 3 lines). No mathematics is written, repaired or re-derived: the statement and the proof are the article's own lines, in the article's own notation, and no retained line is substituted or re-flowed.  No further source-local context is needed: every cross-reference of every retained line resolves inside the two delivered spans.  Nothing was omitted from any retained span, so the package carries no omission record.  No retained line contains a citation, so the wrapper carries no bibliography and no bibliography entry.  No retained line contains a figure environment, an image-inclusion command or drawing code, so no image is delivered and the package declares no asset dependency.  Editorial formatting only, in addition: an AMS \texttt{amsart} preamble that restates the article's own declarations and macros used by the retained lines, namely the environments \texttt{proposition} on separate counters, together with the standard packages those lines need.  The retained environments print in this document as Proposition 1 in order of appearance, whereas the article prints the same items with the numbering of its own full document.  The complete native source of the article as distributed in the arXiv e-print for this exact version is bundled unchanged as \texttt{sources/201878/source0.tex}.
\begin{proposition}
	\label{prop:well-posed}
	Let $(M, g)$ be a Riemannian manifold and $\mathcal{D} = \left\{(t_1, y_1), \ldots, (t_m, y_m)\right\}\subset \RR\times M$ be a finite data set.  Furthermore, let $R\colon TM\to M$ be a retraction and $l_t\colon TM\to TM,\  (x, v)\mapsto (x, tv)$ a vector bundle homomorphism for each $t\in\RR$. We define $V_{y_i}\coloneqq (R\circ l_{t_i})^{-1}(D_{y_i})\subset TM$ and $V\coloneqq \bigcap\limits_{i=1}^m V_{y_i}$. The sets $D_{y_i}$ are chosen such that $R_{y_i} \colon U_{y_i} \to D_{y_i}$ is invertible.
	Then, the mapping defined by 
	\begin{align}
		\label{map: objective function}
		\mathcal{F}_R\colon V\to [0, \infty),\
		(x, v)\mapsto \frac{1}{2} \sum\limits_{i=1}^m \dist_{y_i}^R\left(c_{R, (x, v)}(t_i)\right)^2
	\end{align}
	is well-posed. 
\end{proposition}

\begin{proof}
	Since the retraction-based distance is just defined locally, we have to ensure that the model output $c_{R, (x, v)}(t_i)$ lies in $D_{y_i}$ for each $(x, v)\in V$ and $i=1, \ldots, m$. Therefore, let $i\in\{1, \ldots, m\}$ and $(x, v)\in V$ be arbitrary. According to the definition of $V$ we have especially that $(x,v)\in V_{y_i}$. Thus, the equation $c_{R, (x, v)}(t_i) = R(x, t_iv) = R\left(l_{t_i}(x, v)\right)\in D_{y_i}$ holds for every $i= 1, \ldots, m$ which concludes the proof. 
\end{proof}

\end{document}
